\par\SourceLeaf{142}
Unter solchen Umständen tritt die Wichtigkeit der Function des Vorsitzenden einer berathenden Körperschaft hervor, der die einzelnen Fragepunkte scharf zu sondern, jeden an seinen rechten Platz zu stellen, und in lichtvoller alle Zweideutigkeit ausschliessender Wortfassung zur Berathung und Abstimmung zu bringen hat.

Im vorliegenden Falle war es F., dem als zeitigem Prorector dieses Geschäft oblag. Er setzte unter dem 22. Julius das Rescript, das \textsc{Kritter}’sche Gutachten und das P.’sche P. M. bei sämmtlichen Professoren in Umlauf, mit einer Aufforderung, welche ich in F.’s eignen Worten vollständig hieher setze:
\begin{quote}
Sie möchten sich schriftlich darüber erklären, ob Sie dem ganzen Vorschlage P.’s beitreten, oder über die einzelnen Fragepunkte, nemlich 1) in welchem Maasse die Erhöhung der Pensionen gerecht und rathsam scheine? 2) Ob die Verlängerung der Dauer der Pension nach der Eltern Tode bis zum \(20^{\text{sten}}\) Jahre des jüngsten Kindes gewünscht werde? 3) Die von Königlicher Regierung in Vorschlag gebrachte Erhöhung der jährlichen Beiträge auf 10 Thl. C. M. genehmigt werde, ohne Bedingung; oder unter der Bedingung von Nr. 2.
\end{quote}
\noindent\hspace*{\parindent}Man sieht, dass der letzte Artikel von P.’s P. M. (oben S. [136]) mit keinem Worte erwähnt ist. Ich selbst habe nun zwar keinen Zweifel, dass F. denselben (vielleicht die schwere Wichtigkeit ‘dieses Umstandes’ nicht genug würdigend) als ein schon in seinem Nr. 1 mitenthaltenes Anhängsel betrachtet haben mag: aber eben so wenig zweifle ich, dass diese Nichterwähnung, zumal im Contraste zu der präcis logischen Form, in welcher F. seinen dritten Fragepunkt auftreten lässt, sehr dazu beigetragen hat, jene Hauptfrage für viele, vielleicht für die meisten Votanten in den Hintergrund zu rücken. Ich habe daher geglaubt, diese an sich geringfügigen Nebenumstände hier mitberühren zu müssen, weil dadurch die sonst so auffallende Erscheinung erklärlicher wird, dass von 41 Votanten auch nicht ein einziger sich in eine Discussion über jenen Hauptfragepunkt eingelassen hat (wenn man nicht E.’s Votum S. [139] dafür gelten lassen will); ja dass er von den meisten Votanten gar nicht, und eigentlich nur von zwei Votanten auf ganz unzweideutige Art überhaupt erwähnt ist.

Diese schriftlichen Verhandlungen (vom 22. Julius bis 6. August) sind sehr voluminös, und manche einzelne Abstimmungen sehr ausführlich; allein sie drehen sich fast ausschliesslich um die F.’schen Fragepunkte 2 und 3, welche, besonders der letzte, vielfachen Widerspruch fanden; imgleichen um einige neue im Laufe der Verhandlung eingebrachte Vorschläge, namentlich den einer Aufhebung der bisherigen Freiheit, erst später unter doppelter Nachzahlung der Beiträge in die Witwenkasse eintreten zu können, welcher Vorschlag von einigen Votanten unterstützt, von andern nachdrücklich zurückgewiesen wurde. Ich werde von diesen Abstimmungen nur einige wenige anführen, die mit meinem Gegenstande in näherer Verbindung stehen.

R. erklärt sich überhaupt allen Veränderungen abgeneigt, will aber den Beschlüssen der Mehrheit beitreten. Den P.’schen Art. 3 erwähnt er zwar gar nicht, wohl aber die Principfrage, welche demselben zum Grunde liegt, und in Beziehung auf welche er der von P. bei der frühern Abstimmung aufgestellten Behauptung (oben S. [133] Z. [21]) sehr entschieden entgegentritt.

Er könne, sagt er, sich nicht von der Richtigkeit des Grundsatzes überzeugen, dass bei Erhöhung der jetzigen Witwenpension darauf gesehen werden müsse, dass in der Folge nicht etwa die Nothwendigkeit entstehe, sie zu vermindern, und den künftigen Witwen weniger zu geben, als die jetzigen erhalten. Dieser Grundsatz würde allerdings richtig sein, wenn (wie bei andern Witwenkassen) die Existenz der Kasse bloss auf die Beiträge basirt wäre. Allein, da die Professoren-Witwenkasse ein Beneficium sei, die Beiträge fast für Nichts zu rechnen, und die Theilnahme an dem Beneficium für jeden Professor eine Bedingung seiner Vocation: so gelte, weit natürlicher, der Grundsatz
\begin{quote}
Jede Witwe erhält die möglichst hohe Pension, die der Fonds bei ihrem Lebzeiten verstattet.
\end{quote}
